\documentclass[11pt]{article}
\usepackage[utf8]{inputenc}
\usepackage{amsmath}
\usepackage{geometry}
\geometry{margin=1in}
\title{DEB Equations: V1-morph Feeding}
\author{}
\date{}

\begin{document}
\maketitle
\section{Feeding and V1-morph Implementation}

The equations below follow standard DEB notation: dots indicate rates/fluxes and brackets \([\cdot]\) denote volume-specific quantities.

\subsection{State, environment, and parameters}
\begin{description}
	\item[\(X\)] Food density in the environment (mass or C-moles\,L\(^{-1}\))
	\item[\(X_{\mathrm{in}}\)] Inflow food density, same units as \(X\)
	\item[\(K\)] Half-saturation constant for feeding (same units as \(X\))
	\item[\(k_x\)] Exchange/dilution rate of the environment (time\(^{-1}\))
	\item[\(V\)] Structural volume of the individual (L\(^3\))
	\item[\(f\)] Scaled functional response (dimensionless)
	\item[{\([F_m]\)}] Volume-specific maximum clearance/search rate (time\(^{-1}\)) for V1-morphs
	\item[{\([p_{Am}]\)}] Volume-specific maximum assimilation rate (energy\,L\(^{-3}\)\,time\(^{-1}\))
	\item[\(\kappa_X\)] Assimilation efficiency on food (dimensionless)
	\item[\(\mu_X\)] Food energy density/chemical potential (energy\,mol\(^{-1}\))
	\item[\(\kappa_{\mathrm{bite}}\)] Bite/processing efficiency (dimensionless)
\end{description}

\subsection{Scaled functional response}
\[
	f(X) = \frac{X}{K + X}, \qquad 0 \le f \le 1
\]

\subsection{Environmental dynamics and clearance (chemostat-like, V1-morph)}
The chemostat-like food balance is
\[
	\dot{X} = k_x\,(X_{\mathrm{in}} - X) - F\,X .
\]
For V1-morphs, the individual clearance (feeding) rate scales with structure:
\[
	F = f\,[F_m]\,V \quad \text{(L\(^3\)\,time\(^{-1}\))}.
\]
Substituting into the chemostat food balance yields
\[
	\dot{X} = k_x\,(X_{\mathrm{in}} - X) - f\,[F_m]\,V\,X .
\]

\subsection{Ingestion flux}
Multiplying clearance by environmental density gives the ingestion flux (mass or moles per time):
\[
	\dot{J}_X = F\,X = f\,[F_m]\,V\,X .
\]

\subsection{Assimilation flux (optional, V1-morph)}
In energy units, the assimilation flux can be written as
\[
	\dot{p}_A = f\,[p_{Am}]\,V .
\]
The corresponding molar feeding flux (before assimilation losses) is
\[
	\dot{N}_{\mathrm{feed}} = \frac{\dot{p}_A}{\kappa_X\,\mu_X\,\kappa_{\mathrm{bite}}} = \frac{[p_{Am}]}{\kappa_X\,\mu_X\,\kappa_{\mathrm{bite}}} \frac{X}{K + X}\,V ,
\]
which makes the dependence on \(X\), \(K\), and \(\mu_X\) explicit. Using \(K = \frac{[p_{Am}]}{\kappa_{\mathrm{bite}}\,\kappa_X\,\mu_X\,[F_m]}\) from Section~1 yields
\[
	\dot{N}_{\mathrm{feed}} = \frac{[p_{Am}]\,[F_m]\,X}{[p_{Am}] + \kappa_{\mathrm{bite}}\,\kappa_X\,\mu_X\,[F_m]\,X}\,V ,
\]
so \(\dot{N}_{\mathrm{feed}}\) is expressed entirely in terms of \(\{p_{Am}\}, \{F_m\}, \kappa_X, \mu_X, \kappa_{\mathrm{bite}}\) and the current \(X\).

\end{document}

